\section{Boundary congruences force pure norms and moving Frobenius support}
\label{bc-section}

This section assumes an actual isomorphism of residual Galois modules.
Its arithmetic and density conclusions do not require an eigenform
identification. The final corollary separately uses the boundary orbit
identified in Theorem~\ref{bm-orbit}. No existence of such a congruence,
uniform point-height theorem, or Frey--Mazur statement is assumed.

Put $K=\mathbb Q(r)$, $r^2=-3$, and write
\[
 F=a^4+3a^3b+5a^2b^2+3ab^3+b^4,\qquad
 x=a^2+b^2,\quad y=a+b,
\]
\[
 E_{a,b}:Y^2=X^3+12yX^2+6(3y^2+xr)X,
 \qquad E_0=E_{0,1}.
\]
The local results of Section~\ref{fm-section} give the following actual
seed facts. For coprime positive $a,b$, $F$ is prime to $6$, each
rational prime $q\mid F$ splits in $K$, and both places are
multiplicative with minimal discriminant valuations $m,2m$, where
$m=v_q(F)$. The fixed curve $E_0$ has good reduction above every
$q>3$ and no geometric CM, as proved in the boundary section.

\begin{theorem}[A boundary congruence forces a pure norm]\label{bc-pure}
Let $p>7$ be prime, let $a,b$ be coprime positive integers, and let
$\nu$ be a quadratic character of $G_K$ unramified outside the primes
over $6$, including the trivial character. Suppose
\begin{equation}\label{bc-module}
 E_{a,b}[p]\otimes\overline{\mathbb F}_p
 \simeq (E_0[p]\otimes\overline{\mathbb F}_p)\otimes\nu
 \quad\text{as }G_K\text{-modules}.
\end{equation}
Then $F(a,b)=Q^p$ for an integer $Q>1$. Every prime $q\mid Q$ satisfies
\begin{equation}\label{bc-cutoff}
 q\ge(\sqrt p-1)^2.
\end{equation}
For $q\ne p$, at either place of $K$ above $q$, let $t_q$ be the
trace of $E_0$ over $\mathbb F_q$. Then
\begin{equation}\label{bc-trace-divisor}
 p\mid(q+1)^2-t_q^2.
\end{equation}
In addition, neither $7$ nor $13$ divides $Q$.
\end{theorem}
\begin{proof}
For $q\mid F$ with $q\ne p$, the boundary representation and its
allowed quadratic twist are unramified. At a multiplicative place the
Tate representation modulo $p$ is unramified exactly when the valuation
of its Tate parameter is divisible by $p$. Its valuations here are
$m,2m$, and $p$ is odd. Thus \eqref{bc-module} forces $p\mid m$.

If $p\mid F$, the prime $p$ splits in $K$ and both completions are
$\mathbb Q_p$. The good boundary representation has Serre weight $2$;
the unramified character $\nu$ does not change this inertia condition.
For a multiplicative elliptic curve over $\mathbb Q_p$, the semistable
weight formula of \cite[\S2.9, Proposition~5]{Serre1987Modular} gives
weight $2$ exactly when $p$ divides the negative valuation of $j$.
These negative valuations are $m,2m$, so again $p\mid m$.
Equivalently, the full local representation retains its finite-flat
property under the isomorphism. This argument concerns the full
representations, rather than only local semisimplified traces, and is
unchanged after extending the coefficient field.

Every prime valuation of $F$ is now a multiple of $p$. Since $F\ge13$,
unique factorization gives $F=Q^p$ with $Q>1$.
For a divisor $q\ne p$, the unramified residual Tate representation
has trace $\pm(q+1)$. The additional sign of $\nu$ is absorbed in
this sign, so $t_q\equiv\pm(q+1)\pmod p$, proving
\eqref{bc-trace-divisor}. The Hasse bound makes the corresponding
integer difference nonzero and gives
\[
 p\le |t_q\mp(q+1)|\le q+1+2\sqrt q=(\sqrt q+1)^2.
\]
This is \eqref{bc-cutoff}. If $q=p$, that cutoff holds directly;
no unramified Frobenius assertion at $p$ is used.
The actual depth-one theorem at $13$ excludes $13\mid Q$.
Complete boundary point counts at $7$ give $t_7=\pm4$, so
$p\mid 8^2-4^2=48$ would be necessary if $7\mid Q$. This is
impossible for $p>7$.
\end{proof}

\begin{lemma}[An exact Frobenius count]\label{bc-count}
For an odd prime $p$, put
\[
 C_p=\{A\in\operatorname{GL}_2(\mathbb F_p):
             (\operatorname{tr}A)^2=(\det A+1)^2\}.
\]
Then
\begin{equation}\label{bc-count-formula}
 |C_p|=p(2p^2-p-5),\qquad
 \frac{|C_p|}{|\operatorname{GL}_2(\mathbb F_p)|}
 =\frac{2p^2-p-5}{(p-1)^2(p+1)}.
\end{equation}
\end{lemma}
\begin{proof}
The equation $\operatorname{tr}A=\det A+1$ means that $1$ is an
eigenvalue. For each other eigenvalue $u\in\mathbb F_p^*\setminus\{1\}$,
the class of $\operatorname{diag}(1,u)$ has size $p(p+1)$; there
are $p-2$ choices. With repeated eigenvalue $1$, the identity and the
nontrivial unipotent class together have size $p^2$. Thus the number
of matrices with eigenvalue $1$ is
\[
 (p-2)p(p+1)+p^2=p(p^2-2).
\]
Negation gives the same count for eigenvalue $-1$. Their intersection
is exactly the class of $\operatorname{diag}(1,-1)$, of size $p(p+1)$.
Inclusion--exclusion gives the first formula, and the second follows
from $|\operatorname{GL}_2(\mathbb F_p)|=p(p-1)^2(p+1)$.
\end{proof}

\begin{theorem}[Density of the permitted support]\label{bc-density}
Fix an odd prime $p$ for which the residual image of $E_0$ over $K$
is $\operatorname{GL}_2(\mathbb F_p)$. Among the good prime ideals
away from $p$, let $S_p$ consist of those whose residue-field order
$n$ and trace $t$ satisfy $t^2\equiv(n+1)^2\pmod p$.
Then $S_p$ has natural density
\begin{equation}\label{bc-density-formula}
 \delta_p=\frac{2p^2-p-5}{(p-1)^2(p+1)}
\end{equation}
among the prime ideals of $K$, ordered by norm. The corresponding
permitted split rational primes have natural density $\delta_p/2$
among all rational primes. The full-image hypothesis holds for every
sufficiently large $p$, with no effective threshold asserted here;
along those primes, $\delta_p\sim2/p$.
\end{theorem}
\begin{proof}
In $K(E_0[p])/K$, Frobenius has trace $t$ and determinant $n$ modulo
$p$. The condition defining $C_p$ is conjugacy invariant. Chebotarev
natural density, in the form of \cite[Definition~8.30 and
Theorem~8.31]{MilneANT2020}, and Lemma~\ref{bc-count} give
\eqref{bc-density-formula}.

At a split rational prime, the traces at the two places have the same
square. Indeed, the explicit two-isogeny relates $E_0$ to the
$(-2)$-twist of its conjugate, so these traces differ by the quadratic
sign $(-2/q)$. Hence both places are permitted or neither is.
Prime ideals above inert rational primes counted up to norm $X$ make
only an $O(\sqrt X)$ contribution. Each permitted split rational
prime contributes two prime ideals, giving density $\delta_p/2$
among rational primes. The split primes themselves have density $1/2$.

Finally $E_0$ has no geometric CM. Serre's open-image theorem gives
full mod-$p$ image over the fixed number field for all but finitely
many $p$; see \cite[pp.~259--260, assertions (3), (6), (7)]{Serre1972OpenImage}.
The displayed expression for $\delta_p$ has the stated asymptotic.
\end{proof}

These density assertions concern each fixed $p$. They provide no
Chebotarev error bound uniform in $p$ and the norm cutoff. For every
fixed sufficiently large $p$, the permitted support still has positive
density and is infinite. The allowed twists $\nu$ only change traces
by sign away from $6p$, so this support set is unchanged; no image
threshold uniform over varying twists is required. Under
\eqref{bc-module}, every prime divisor other than possibly $p$ belongs
to this permitted split support and also satisfies \eqref{bc-cutoff}.
Neither restriction alone excludes all prime factors of a polynomial
value from lying in that support.

\begin{corollary}[All embeddings of the boundary orbit]\label{bc-orbit}
Suppose the descended residual representation of a coprime positive
seed, at a prime $p>7$, is congruent to any coefficient embedding of
the level-$576$ boundary orbit identified in Theorem~\ref{bm-orbit}.
Then $F(a,b)$ is a pure $p$-th power and satisfies the support
conditions of Theorem~\ref{bc-pure}. In particular, in the canonical
decomposition $F=VQ^p$ with $V$ $p$-power-free, that orbit cannot
occur when $V>1$.
\end{corollary}
\begin{proof}
The specified order-four character in the descent is $\chi$.
Compatibility and the algebraic Frobenius identities of the boundary
identification show that restriction to $G_K$ compares
$E_{a,b}[p]\otimes\chi$ with $E_0[p]\otimes\chi^\sigma$.
This comparison uses the compatible representations, not finitely
many coefficient values. The checked large-image input gives absolute
irreducibility for the actual representation, so the resulting
semisimplified isomorphism is the full module isomorphism needed above.
An embedding either fixes or inverts the values of $\chi$, and hence
\[
 \chi^\sigma/\chi\in\{1,\chi^{-2}\}
       =\{1,\psi_3|_{G_K}\}.
\]
Both possibilities are quadratic and unramified outside $6$.
Cancelling the specified $\chi$ therefore gives
\eqref{bc-module}, possibly with its nontrivial quadratic twist.
Theorem~\ref{bc-pure} applies. This does not exclude the pure branch.
\end{proof}

\paragraph{Verification and remaining scope.}
The ordinary proofs have independent reviews. Exact standard-library
replay exhausts all matrices for $p=3,5,7,11,13,17,19$ and all points
of the eight boundary reductions over $q=7,13,19,31$. The general
matrix formula is proved above; the finite checks supplement it.
The local trace divisibility test always retains the condition $q\ne p$.
The orbit corollary has the additional modular-identification dependency
stated in its hypothesis. No Lean formalization of the Tate, Serre,
Chebotarev or modular comparison theorems is asserted by these checks.
